\section{Conditional Continuation Comparisons and Decision-Class Accuracy}
\subsection{A direct conditional payoff statistic}
\label{app:r16menustatistic}

For one fixed query and two already selected actions $a,b$, let
$\overline F(x)=\{F(x;\xi)+F(x;-\xi)\}/2$, where a new base array
$\xi$ is independent of fitting and all other confirmation pairs. Put
$\gamma=2$ times the stored terminal-dispersion penalty and define
\[
 q_0^{ab}=-\frac\gamma2
       \{\|\Pi x_a\|_d^2-\|\Pi x_b\|_d^2\},\qquad
 R^{ab}=\overline F(x_a)-\overline F(x_b)-q_0^{ab}.
\]
The current reward and $q_0^{ab}$ are known conditional on this query
and its actions. Thus
\begin{equation}
 Q_\theta(y,a)-Q_\theta(y,b)
       =r_\theta(y,a)-r_\theta(y,b)+q_0^{ab}+\E R^{ab}.
 \label{eq:r16menucentered}
\end{equation}
The implementation may center at an enclosed binary64 representative of
$q_0^{ab}$. Its exact discrepancy is included in the range; the same
representative is added back to the known part of the payoff.

Let $N$ denote the number of dates, $n=N-1$, and let $\nu_I,\nu_C$
denote the stored idiosyncratic and common innovation amplitudes per
date. Set
\[
 H=(1+h\beta)^n,\qquad
 P=\beta_\psi\sum_{j=1}^{N-1}B_j(1+h\beta)^{j-1},\qquad
 D=\|x_a-x_b\|_d,
\]
and $s=\max\{\|\Pi x_a\|_d,\|\Pi x_b\|_d\}$. The following
coefficients contain no learned parameter:
\begin{align}
 C&=P+\gamma H\kappa nh
       +\gamma(H-1)\{s+\nu_I\sqrt{N(1-1/d)}\},\notag\\
 A&=\gamma(H-1)\nu_I\sqrt{N/d}.
 \label{eq:r16menuconstants}
\end{align}

\begin{proposition}[Query-conditional antithetic range]
\label{prop:r16menurange}
For one independent antithetic pair,
\[
 \Pr\{|R^{ab}|>D(C+Av)\}\leq e^{-v^2/2},\qquad v\geq0.
\]
Consequently, clipping at the prospectively fixed $v_*>0$ admits
the expectation allowance $DAe^{-v_*^2/2}/v_*$. These statements
hold conditional on the selected actions, so their geometry may enter
the range before the independent pair is generated.
\end{proposition}

\begin{proof}
Apply the pathwise mean-value identity along the straight segment from
$x_b$ to $x_a$. The reference adjoint has norm growth at most $H$;
its accumulated production component has norm at most $P$. Subtracting
the known terminal derivative $-\gamma\Pi x$ leaves the production
adjoint, a bounded accumulated-production term, and
$(J^\top-I)$ times the centered starting state and centered future
innovation. The unmultiplied terminal innovation cancels exactly in an
antithetic pair. The starting state is fixed. The norm of the projected
future innovation has mean at most $\nu_I\sqrt{N(1-1/d)}$ and,
as a function of its standard Gaussian array, normalized Lipschitz
constant at most $\nu_I\sqrt{N/d}$. Gaussian concentration gives
one event, uniform over the entire straight segment. Integration proves
the range, and integration of its tail proves the expectation allowance.
There is no prefix-state event because the query is conditioned upon.
\end{proof}

One observation is one antithetic pair, not its two dependent paths.
Different query-pair cells use independent arrays. Equal pair counts
at every query and at every declared training stream identify their
uniform finite average. For each contrast, the report pools only its
independent pair observations, clips each query at its proved range,
uses the maximum range in the independent nonidentical empirical-Bernstein
bound, and averages numerical and tail allowances. Dependence across
methods is retained by forming the direct contrast before inference.
The primary family comprises 216 two-sided events and uses probability
budget $0.018$. The separate scalar-accuracy family below uses $0.002$.
Their union remains within the fixed total $0.02$ budget.

\paragraph*{Arithmetic and Gaussian clipping.}
The numerical verifier uses the preserved protected polynomial for
$\tanh$, outward coefficient operations, and forward state-error
recursions with multiplier $1+h\beta$. Fixed dyadic spectral proposals
for the two declared coupling matrices are each certified by a fresh
outward LDL calculation. Postdecision representatives are interval
midpoints, so reconstructing their error accounts does not depend on
the last bits of a numerical singular-value or matrix-product routine.
A deterministic state cap
is computed from the fixed query, bounded drift and reference actions,
and the executed Gaussian cap. Production error is bounded by
$\beta_\psi$ times the state error plus its protected-function and
dot-product allowances. A terminal state error $e$ at a cap $M$
contributes at most $(\gamma/2)e(2M+e)$, before the separately charged
scalar reductions. The source records every component of this account.

The finite economy retains untruncated Gaussian innovations. For cap $z$,
\[
 \E\{G-\operatorname{clip}_{z}(G)\}^2
       \leq 4\phi(z)/z^3.
\]
This follows by writing the two-sided tail integral as
$2\phi(z)\int_0^\infty u^2e^{-zu-u^2/2}du$ and dropping
$e^{-u^2/2}$. Multiplication by $\nu_I^2+\nu_C^2$ and the
same state-error recursion gives an $L^2$ clipping allowance. The
production Lipschitz bound and conditional terminal Cauchy--Schwarz
then transfer the implemented expectation to the original finite
Gaussian economy. The clipping and arithmetic accounts are distinct.

\subsection{Accuracy over an entire scalar intervention class}
\label{app:r16menuaccuracy}

A separate comparison supplies an absolute reference on a restricted
decision class while retaining the high-dimensional state. At one fixed
state and current task, set $a(s)=a_L+s(a_R-a_L)$, $0\leq s\leq1$.
The endpoints are the common first-period reference consumption plus
or minus the fixed radius, intersected with the original task box.
All subsequent actions follow the same reference. Write
$q(s)=Q_\theta\{y,a(s)\}$ and $r=a_R-a_L$.

For the first and second variations of a future state with respect to
$s$, start with $V_1=h\|r\|_d$ and $W_1=0$. With the global
scalar and vector Hessian coefficients $H_m,H_x$ from the paired
transfer, recursively set
\[
 V_{j+1}=(1+h\beta)V_j,\qquad
 W_{j+1}=(1+h\beta)W_j+hH_xV_j^2.
\]
Let
\[
 S_N=\max\{\|\Pi x_{a_L}\|_d,\|\Pi x_{a_R}\|_d\}
          +\kappa(N-1)h+\nu_I\sqrt{N(1-1/d)}.
\]
Then an upper continuation curvature is
\[
 \Lambda=\sum_{j=1}^{N-1}B_j(H_mV_j^2+\beta_\psi W_j)
                                         +\gamma S_NW_N.
\]
For $u_i=\max(a_{L,i},a_{R,i})$, current utility weight $w$ and
adjustment coefficient $\eta$, put
\[
 \mu=A_0\left\{w\overline{r_i^2/u_i^2}+\eta\bar r^2\right\}-\Lambda.
\]

\begin{proposition}[Scalar Bellman curvature and a protected optimality gap]
\label{prop:r16menugap}
The preceding coefficients imply $q''(s)\leq-\mu$ throughout
$[0,1]$. If $[g_L,g_U]$ is a valid interval for $q'(s_0)$, then
\[
 \max_{0\leq s\leq1}q(s)-q(s_0)
 \leq\max_{\substack{g\in\{g_L,g_U\}\\0\leq s\leq1}}
           \{g(s-s_0)-\mu(s-s_0)^2/2\}.
\]
The maximum is a scalar quadratic calculation. It remains valid when
$\mu\leq0$; in that case it does not assert strong concavity.
\end{proposition}

\begin{proof}
Differentiating the finite transition recursion gives the displayed
variation bounds. The production chain rule contributes at most
$H_mV_j^2+\beta_\psi W_j$. The terminal second derivative is
$-\gamma\{\|\Pi v_N\|_d^2+\langle\Pi Y_N,\Pi w_N\rangle_d\}$.
Dropping its negative squared term and using the deterministic bound
on $w_N$ and the moment bound on $Y_N$ gives $\gamma S_NW_N$.
Subtracting the minimum curvature of current utility proves the
first claim. Taylor's integral remainder gives the quadratic support
at $s_0$. Maximizing that support over the gradient interval and
feasible scalar displacement proves the second.
\end{proof}

The gradient interval is obtained without identifying an automatic
derivative with an exact expectation. On a fixed neighboring interval
$[s_-,s_+]$ containing $s_0$, a fresh paired-payoff calculation bounds
$\{q(s_+)-q(s_-)\}/(s_+-s_-)$. A global absolute second-derivative
bound $L$ adds the deterministic allowance
\[
 \frac{L\{(s_0-s_-)^2+(s_+-s_0)^2\}}{2(s_+-s_-)}.
\]
This is the integral of the derivative's Lipschitz bound; it also
covers one-sided differences at a boundary. The declared scalar
candidate is selected from the fitted continuation before this bank.
Binary64 construction of the affine actions receives its own payoff
allowance. All 128 scalar events use 65,536 independent antithetic
pairs each and the fixed probability allocation. The resulting
reference covers every scalar action in the interval, including
unevaluated ones. Its scope is the specified finite decision class;
it is not a high-dimensional full-adapted-control or continuous-time
near-optimality assertion.

\subsection{Scope of the scalar accuracy reference}
The protected quadratic upper bound controls the best action over the entire declared scalar intervention segment, not merely a finite grid of evaluated actions. It is therefore a distinct absolute accuracy object from the broad adapted-control bound. It neither covers all high-dimensional current actions nor changes the original continuous-time control problem into a finite scalar problem. The stored finite-model coefficients, horizon, task law, scalar segment, and subsequent reference policy are part of its definition.

The interval law, numerical transfer, and whole-segment curvature are proved before their statistical use. A manufactured check of these identities verifies the calculation, not a successful economic outcome for a learned candidate. When a numerical scalar bound is reported, its event, selected scalar candidate, and comparison class must accompany it. No numerical value is inferred from the positivity of a curvature constant alone.

\subsection{Decision value with binding task constraints}
\label{sec:r16activefacevalue}
Temporary action constraints need not rule out a strict decision-value result. The interior assumption in Proposition~\ref{prop:r16decisionvalue} can be replaced by a stable-face condition. The relevant prediction norm then charges errors only in the directions in which each economic task can respond.

Retain the quadratic objectives and the common observation $z=q+\varepsilon$ of Section~\ref{sec:r16decisionvalue}, with $\E\varepsilon=0$ and covariance $\Sigma$. For each task, fix a nonempty face $F_j$ of its feasible polytope, a point $a_j^0\in F_j$, and a full-column-rank matrix $L_j$ spanning the parallel space of $F_j$. Put
\[
 R_j=L_j(L_j^\top H_jL_j)^{-1}L_j^\top,
 \qquad W_F=\sum_{j=1}^K\omega_jh_j^2R_j.
\]
For a singleton face set $R_j=0$. Assume that $W_F$ is positive definite. Individual faces may have binding constraints; the last condition requires their decision-relevant tangent spaces jointly to span the continuation coordinates.

Let $D$ be a fixed scalar-gradient dictionary, independent of $\varepsilon$, and let $P_F$ project orthogonally onto the columns of $W_F^{1/2}D$. Define
\[
 \widehat q_R=z,\qquad
 \widehat q_C=W_F^{-1/2}P_FW_F^{1/2}z,
 \qquad y_F=W_F^{1/2}q,
 \quad\Omega_F=W_F^{1/2}\Sigma W_F^{1/2}.
\]
The stable-face condition is that, for $q$, $\widehat q_R$, and $\widehat q_C$, the unique constrained maximizer lies in the relative interior of the same deterministic face $F_j$, almost surely, for every task. This is a condition on the original observation law; it does not condition an otherwise unrestricted noise distribution on a face-selection event. It can be verified by the inactive slacks and the normal optimality inequalities on the prescribed noise support.

\begin{proposition}[Continuation compression on stable active faces]
\label{prop:r16activefacevalue}
Under these conditions, the expected weighted payoff difference between the constrained critic and Raw actions equals
\begin{equation}
 \frac12\left[\operatorname{tr}\{(I-P_F)\Omega_F\}
                  -\|(I-P_F)y_F\|^2\right].
 \label{eq:r16activefacevalue}
\end{equation}
Consequently, compression strictly improves expected economic payoff whenever the removed decision-relevant variance exceeds the weighted approximation cost. In particular, this holds for a correctly specified dictionary that removes strictly positive decision-relevant variance, even though every task may have binding constraints.
\end{proposition}
\begin{proof}
Tangential optimality on $F_j$ gives, for each of the three continuation vectors $v$,
\[
 a_j(v)=a_j^0+R_j(b_j-H_ja_j^0-h_jv).
\]
Thus $a_j(v)-a_j(q)=-h_jR_j(v-q)$. The gradient of the true objective at $a_j(q)$ is orthogonal to this difference. Expanding the quadratic objective, and using $R_jH_jR_j=R_j$, yields
\[
 Q_j(a_j(q);q)-Q_j(a_j(v);q)
       =\tfrac12 h_j^2(v-q)^\top R_j(v-q).
\]
Summing the weighted losses gives Raw expected loss
\[
 \tfrac12\operatorname{tr}(\Omega_F).
\]
For the critic, the weighted error is
\[
 -(I-P_F)y_F+P_FW_F^{1/2}\varepsilon.
\]
Its expected squared norm is
\[
 \|(I-P_F)y_F\|^2+\operatorname{tr}(P_F\Omega_F).
\]
Subtracting the two expected losses proves the formula and its strictness assertions.
\end{proof}

The result gives a positive boundary-aware mechanism, not a universal ranking of methods. The faces and dictionary must satisfy the stated conditions for the full observation law. If fitted actions switch faces, the displayed exact identity cannot be obtained merely by selecting the observations with a favorable face. If $W_F$ is singular, the present statement does not invoke an ordinary inverse; a lower-dimensional formulation requires a separately specified decision space. The nonlinear menu is evaluated by its own paired payoff intervals, rather than by asserting that its fitted features or active sets satisfy this specialization. Identical projected estimators, whether called Raw or NBO, still have identical decisions.

